% !TeX root = ../main.tex
%%% ---------------------------------------------------------------------------
%%% 方法
%%%
%%% 这一节演示了公式、定理、算法的规范写法，改写时请保留这些约定：
%%%   · 变量斜体、单位正体、函数名正体（GB/T 3101，见 ../../STANDARDS.md §1.3）
%%%   · 微分写 \dd x，自然常数写 \e，虚数单位写 \ii
%%%   · 说明性下标写 \subtext{max}，不要写 _{max}
%%%   · 矢量矩阵用 \vect{} \mat{}，不要用 \mathbf 或 \boldsymbol
%%%   · 每个公式都给 \label，引用时用 \cref
%%% ---------------------------------------------------------------------------
\section{方法}
\label{sec:method}

\subsection{总体结构}
\label{sec:method:overview}

设骨干网络输出的层级特征为 $\setof{\featmap_l}_{l=2}^{5}$，其中 $\featmap_l \in
\R^{C \times H_l \times W_l}$，$H_l = H/2^{\,l}$。\ourmethod{} 在自顶向下
路径的每一次融合处替换原有的逐元素相加，代之以\cref{sec:method:gate} 定义的
门控算子。整体结构如\cref{fig:architecture} 所示。

\subsection{跨层注意力门控}
\label{sec:method:gate}

记上层特征经上采样对齐后为 $\tilde{\featmap}_{l+1}$。门控权重
$\mat{G}_l \in [0,1]^{H_l \times W_l}$ 由两层特征的通道统计量共同决定：

\begin{equation}
  \label{eq:gate}
  \mat{G}_l = \sigma\!\left(
    \mat{W}_2 \relu\!\left(
      \mat{W}_1 \left[ \operatorname{GAP}(\featmap_l) \,\Vert\,
                       \operatorname{GAP}(\tilde{\featmap}_{l+1}) \right]
    \right)
  \right),
\end{equation}

式中 $\sigma(\cdot)$ 为 Sigmoid 函数，$\operatorname{GAP}$ 为全局平均池化，
$\Vert$ 表示通道维拼接，$\mat{W}_1 \in \R^{C/r \times 2C}$ 与
$\mat{W}_2 \in \R^{C \times C/r}$ 为可学习参数，压缩比取 $r = 16$。

融合结果为

\begin{equation}
  \label{eq:fusion}
  \featmap_l' = \mat{G}_l \odot \featmap_l
              + (\mathbf{1} - \mat{G}_l) \odot \mathcal{D}(\tilde{\featmap}_{l+1}),
\end{equation}

其中 $\odot$ 为逐元素乘，$\mathcal{D}(\cdot)$ 为\cref{sec:method:align}
中的可变形对齐算子。由\cref{eq:fusion} 可见，当 $\mat{G}_l \to \mathbf{1}$ 时
融合退化为只保留浅层特征，当 $\mat{G}_l \to \mathbf{0}$ 时退化为标准
特征金字塔，因此本文方法是后者的严格推广。

\subsection{可变形对齐}
\label{sec:method:align}

上采样引入的偏移量 $\Delta \vect{p}$ 由 $3 \times 3$ 卷积从拼接特征回归：

\begin{equation}
  \label{eq:offset}
  \mathcal{D}(\tilde{\featmap})(\vect{p})
    = \sum_{k=1}^{9} w_k \cdot
      \tilde{\featmap}\!\left(\vect{p} + \vect{p}_k + \Delta \vect{p}_k\right),
\end{equation}

$\vect{p}_k$ 为标准卷积核的采样网格，$w_k$ 为核权重。

\subsection{理论分析}
\label{sec:method:theory}

下面说明门控融合不会放大特征的方差，这是训练稳定性的必要条件。

\begin{assumption}
  \label{assum:indep}
  $\featmap_l$ 与 $\mathcal{D}(\tilde{\featmap}_{l+1})$ 在各空间位置上
  互不相关，且方差分别为 $\sigma_1^2$ 与 $\sigma_2^2$。
\end{assumption}

\begin{theorem}[方差不增]
  \label{thm:variance}
  在\cref{assum:indep} 下，对任意 $\mat{G}_l \in [0,1]^{H_l \times W_l}$，
  由\cref{eq:fusion} 给出的融合特征满足
  \begin{equation}
    \label{eq:var-bound}
    \Var\!\left[\featmap_l'\right] \le \max\setof{\sigma_1^2, \sigma_2^2}.
  \end{equation}
\end{theorem}

\begin{proof}
  记 $g \in [0,1]$ 为 $\mat{G}_l$ 在某一位置的取值。由独立性，
  \begin{equation}
    \label{eq:var-expand}
    \Var\!\left[\featmap_l'\right]
      = g^2 \sigma_1^2 + (1-g)^2 \sigma_2^2 .
  \end{equation}
  右端是关于 $g$ 的二次函数，在 $[0,1]$ 上的最大值只可能在端点取得，
  分别为 $\sigma_1^2$ 与 $\sigma_2^2$，即得\cref{eq:var-bound}。
\end{proof}

\begin{remark}
  \label{rmk:relax}
  \cref{assum:indep} 的独立性在实践中并不严格成立，但实验中未观察到
  训练发散，说明该界在相关性较弱时仍具指导意义。
\end{remark}

\subsection{训练流程}
\label{sec:method:train}

完整的训练流程见\cref{alg:training}。

\begin{algorithm}[htbp]
  \caption{\ourmethod{} 的训练流程}
  \label{alg:training}
  \begin{algorithmic}[1]
    \Require 训练集 $\mathcal{S}$，骨干网络参数 $\params\subtext{bb}$，
             学习率 $\eta$，总轮数 $T$
    \Ensure  收敛后的模型参数 $\params^{*}$
    \State 用 ImageNet 预训练权重初始化 $\params\subtext{bb}$，其余参数用
           Kaiming 初始化
    \For{$t \gets 1$ \textbf{to} $T$}
      \ForAll{小批量 $\mathcal{B} \subset \mathcal{S}$}
        \State 前向计算层级特征 $\setof{\featmap_l}$
        \For{$l \gets 4$ \textbf{down to} $2$}
          \State 按\cref{eq:gate} 计算门控权重 $\mat{G}_l$
          \State 按\cref{eq:fusion} 得到融合特征 $\featmap_l'$
        \EndFor
        \State 计算总损失 $\loss = \loss\subtext{cls} + \lambda \loss\subtext{reg}$
        \State $\params \gets \params - \eta \nabla_{\params} \loss$
      \EndFor
      \State 按余弦策略更新 $\eta$
    \EndFor
    \State \Return $\params^{*} \gets \argmin_{\params} \loss$
  \end{algorithmic}
\end{algorithm}
